\addtotoc{Abstract}
\begin{abstract}
	\addtocontents{toc}{}

	\hspace{1cm} Autonomous penetration testing --- the use of LLM-orchestrated agents to discover and exploit vulnerabilities in a target system without step-by-step human direction --- has advanced rapidly since 2023, yet two quantitatively documented failure modes recur across independently developed systems: CVE-Bench finds that insufficient exploration accounts for 37.5\%--80.0\% of all unsuccessful exploitation attempts, with the best-performing agent achieving only a 13\% pass@5 success rate on 40 critical-severity real-world web vulnerabilities; and PentestEval finds that Attack Decision-Making --- the stage requiring prerequisite reasoning between discovered weaknesses --- is the weakest planning stage across nine evaluated LLMs, averaging a Spearman correlation of 0.25 against expert judgement, with a ground-truth-injection ablation showing that expert-quality ADM alone raises end-to-end success from 0.31 to 0.67. No system in the reviewed literature combines open-ended discovery of an unfamiliar attack surface with a dynamically constructed, prerequisite-aware model of vulnerability dependencies, and no autonomous system has been evaluated across more than one independently benchmarked attack-surface family. This inception report documents the first stage of an investigation into closing this compound gap through RedGrid: an LLM-orchestrated multi-agent system built around a four-layer orchestration hierarchy (Mission Planner, Team Manager, Specialist Agents, and Execution/Validation Layer), a Dual-Layer World Model that separates confirmed evidence from inferred hypotheses, a Vulnerability Dependency Graph (VDG) with Upper Confidence Bound (UCB) guided node selection for prerequisite-aware exploration, and a Cross-Mission Memory module for selective strategy reuse across engagements. The evaluation plan targets four existing oracle-backed benchmarks spanning three attack-surface families: CVE-Bench and PentestEval (web applications), PrediQL (GraphQL APIs), and Incalmo MHBench (multi-host Active Directory environments). At this stage no results are available; the report establishes the research direction, motivates it against the quantitative evidence from the reviewed literature, and defines the proposed architecture and evaluation plan for the work ahead.

\end{abstract}

\clearpage
\setstretch{1.5}
